\section{An angular-volume criterion for the critical tail}
The dimension thresholds above come from the geometry of the maximizing
correlation directions. The following elementary criterion separates that
geometry from the particular parity calculation. It is a Laplace-method
calculation, not a new general theory of heavy-tailed condensation.

Let $U\sim N(0,I_{d+1})$, $a\sim N(0,1)$ independently, and now let
$G:\mathbb R^{d+1}\to\mathbb R$ be any continuous positively
homogeneous function of degree one. Write
$g(\omega)=G(\omega)$ for $\omega\in S^d$ and assume
$\rho=\max_{S^d}|g|>0$. Surface measure below is unnormalized.
Define the angular near-maximum volume
\[
 V(u)=\operatorname{area}\left\{\omega\in S^d:
        |g(\omega)|>0,\quad |g(\omega)|^{-1}-\rho^{-1}\le u\right\}.
\]
Suppose, for some $L>0$ and $\kappa>0$,
\[
                         V(u)\sim L u^\kappa\qquad(u\downarrow0).
 \tag{25}
\]
The zero set of $g$ is allowed to have positive surface measure.

\begin{theorem}[Angular volume determines endpoint moments]
Put $Z=aG(U)$ and $t_*=1/\rho$, and set
$\alpha=\kappa-(d+1)/2$. For each integer $k\ge0$,
\[
        M^{(k)}(t_*-)<\infty\quad\Longleftrightarrow\quad\alpha>k.
 \tag{26}
\]
If $\alpha>0$, the critical law $Q_*=e^{t_*Z}P_0/M(t_*)$ exists,
and its density on the positive half-line satisfies
\[
 f_*(z)\sim A z^{-\alpha-1},\qquad
 A=\frac{(2\pi)^{-(d+2)/2}}{M(t_*)}
       \sqrt{\frac\pi2}\,
       \rho^{-(d+1)/2}L\Gamma(\kappa+1).
 \tag{27}
\]
Consequently $Q_*(Z>z)\sim (A/\alpha)z^{-\alpha}$.
If $\alpha>1$, its mean $m_c$ is finite and the scalar entropy
problem above $m_c$ has the same linear, unattained infimum as in
Theorem~\ref{thm:noattain}. This last assertion concerns only the scalar
constraint, without an additional symmetry hypothesis on output vectors.
\end{theorem}
\begin{proof}
After integrating $a$, polar coordinates express the endpoint integral,
up to a positive constant, as
\[
 \int_0^\infty R^d\int_{S^d}
       \exp\{-R^2[1-g(\omega)^2/\rho^2]/2\}\,d\omega\,dR.
\]
Near the maximum,
$1-g^2/\rho^2\sim2\rho(|g|^{-1}-\rho^{-1})$.
Integration by parts in the angular sublevel volume (25) gives an
angular integral asymptotic to a positive constant times $R^{-2\kappa}$.
The part away from the maximum is exponentially small. Inserting
$G(U)^{2k}=R^{2k}g(\omega)^{2k}$ changes this asymptotic by a
positive constant times $R^{2k}$. The radial tail therefore converges
exactly when $d+2k-2\kappa<-1$. As in the proof of
Theorem~\ref{thm:boundary}, differentiating after integrating $a$
produces a polynomial in $G^2$ with positive leading coefficient and
nonnegative other coefficients at the positive endpoint. This proves
(26), including logarithmic divergence at equality.

For the density, integrate $a=z/G(U)$, and then use polar coordinates.
With $q=|g(\omega)|>0$ this gives
\[
 f_*(z)=\frac{(2\pi)^{-(d+2)/2}}{M(t_*)}e^{z/\rho}
 \int_{S^d}\frac1q\int_0^\infty
 R^{d-1}\exp\left\{-\frac{R^2}2-\frac{z^2}{2q^2R^2}\right\}
                                           dR\,d\omega. \tag{28}
\]
Directions with $q=0$ contribute only to the atom at zero.
The elementary radial saddle estimate is
\[
 \int_0^\infty R^{d-1}e^{-R^2/2-z^2/(2q^2R^2)}dR
 \sim\sqrt{\frac\pi2}\,(z/q)^{(d-1)/2}e^{-z/q}.
 \tag{29}
\]
It is uniform for $q$ in a fixed neighborhood of $\rho$.
For example set $R=\sqrt{z/q}\,y$; the exponent has its unique
minimum at $y=1$, with value $z/q$ and second derivative $4z/q$.
Taylor expansion on a shrinking neighborhood of one, followed by
Gaussian integration, gives (29). The complement has a strict
exponential gap. This proof also gives an upper bound by a constant
multiple of its right side whenever $z/q\ge1$.

The angular region $q\le\rho-\epsilon$ contributes exponentially
less than the main term: the last bound, with its fixed power of
$q^{-1}$, is dominated by $C e^{-c_\epsilon z}$ after factoring
out $e^{-z/\rho}$. In the remaining region the factor
$q^{-(d+1)/2}$ tends to $\rho^{-(d+1)/2}$ at the maximizing set.
Finally (25) and integration by parts give
\[
 \int_{S^d}e^{-z(q^{-1}-\rho^{-1})}d\omega
       \sim L\Gamma(\kappa+1)z^{-\kappa}.
\]
Substitution in (28)--(29) proves (27). Integrating the density gives
the tail. Readout symmetry again makes the negative density
$e^{-2t_*z}f_*(z)$.

For the entropy assertion, its lower bound and nonattainment follow
from the critical exponential-tilt inequality exactly as before.
For a matching sequence let $B_R=\{R<Z\le R+1\}$ and
$\Pi_R=Q_*(\,\cdot\mid B_R)$. Equation (27) gives
$Q_*(B_R)\sim A R^{-\alpha-1}$, whence
\[
 \operatorname{KL}(\Pi_R\|P_0)
 =t_*\E_{\Pi_R}Z-\log M(t_*)-\log Q_*(B_R)
                  =t_*R+O(\log R).
\]
Its mean lies in $[R,R+1]$. Mixing it with $Q_*$ in the unique
proportion giving mean $s$ and using convexity of relative entropy
matches $t_*s-\log M(t_*)$ in the limit. Equality would require
$Q=Q_*$, whose mean is too small.
\end{proof}

Suppose, more specifically, that the maximizing set is a compact smooth
$m$-dimensional submanifold of the $d$-dimensional sphere, and the
correlation deficit is comparable to the $\beta$th power of its normal
distance. Tubular coordinates then give
$V(u)\asymp u^{(d-m)/\beta}$. Exact asymptotics require the corresponding
leading normal-volume limit, but the two-sided comparison already
suffices for the endpoint integrability test. Thus the exponent is
\[
                  \alpha=\frac{d-m}{\beta}-\frac{d+1}{2}. \tag{30}
\]
An isolated sharp corner has $m=0,\beta=1$ and gives
$\alpha=(d-1)/2$, as for ReLU parity. An isolated smooth nondegenerate
quadratic maximum has $m=0,\beta=2$ and gives $\alpha=-1/2$:
the endpoint normalizer then diverges, even though the maximizing
directions themselves form a finite set. Linear $G(U)$ supplies an
explicit example of the latter case.

For ReLU parity, the corner identity gives
\[
 L\Gamma(d+1)=\rho^d\sum_s\int_{v_s^\perp}e^{-h_s(t)/\rho}\,dt.
\]
Indeed its angular gap is $h_s(u)/\rho^2+o(\|u\|)$, and the
volume of a norm ball scales with the $d$th power of its radius.
Formula (27) therefore reduces exactly to $A=\alpha C$ in (10)--(11).
For leaky ReLU with fixed $0\le\lambda<1$,
$\phi_\lambda(t)=(1-\lambda)\relu(t)+\lambda t$, the parity
statistic for even $r\ge2$ is $(1-\lambda)G(U)$: its affine part
has zero parity coefficient. It has the same exponent and moment
thresholds, with the critical tilt rescaled by $(1-\lambda)^{-1}$.
